(a) $M_2 = \MbP{0}{2}$ (Uhrzeigersinn) $<$ $M_3 = \McP{0}{2}$ (Gegenuhrzeigersinn) $<$ $M_1 = \MaP{0}{2}$ (Gegenuhrzeigersinn), (b) beide erhalten $F_3, F_1, F_2$; Lea vergisst den Hebelarm, Tim misst bis zum Angriffspunkt statt zur Wirkungslinie ($F_2\cdot r_2 = \MwP{0}{2}$ statt $\MbP{0}{2}$)
